\documentclass{amsart}
% For dealing with references we use the comment environment
\usepackage{verbatim}
\newenvironment{reference}{\comment}{\endcomment}
%\newenvironment{reference}{}{}
\newenvironment{slogan}{\comment}{\endcomment}
\newenvironment{history}{\comment}{\endcomment}

% For commutative diagrams we use Xy-pic
\usepackage[all]{xy}

% We use 2cell for 2-commutative diagrams.
\xyoption{2cell}
\UseAllTwocells

% We use multicol for the list of chapters between chapters
\usepackage{multicol}

% This is generally recommended for better output
\usepackage{lmodern}
\usepackage[T1]{fontenc}

% For cross-file-references

% Package for hypertext links:
\usepackage{hyperref}

% For any local file, say "hello.tex" you want to link to please

% Theorem environments.
%
\theoremstyle{plain}
\newtheorem{theorem}[subsection]{Theorem}
\newtheorem{proposition}[subsection]{Proposition}
\newtheorem{lemma}[subsection]{Lemma}

\theoremstyle{definition}
\newtheorem{definition}[subsection]{Definition}
\newtheorem{example}[subsection]{Example}
\newtheorem{exercise}[subsection]{Exercise}
\newtheorem{situation}[subsection]{Situation}

\theoremstyle{remark}
\newtheorem{remark}[subsection]{Remark}
\newtheorem{remarks}[subsection]{Remarks}

\numberwithin{equation}{subsection}

% Macros
%
\def\lim{\mathop{\mathrm{lim}}\nolimits}
\def\colim{\mathop{\mathrm{colim}}\nolimits}
\def\Spec{\mathop{\mathrm{Spec}}}
\def\Hom{\mathop{\mathrm{Hom}}\nolimits}
\def\Ext{\mathop{\mathrm{Ext}}\nolimits}
\def\SheafHom{\mathop{\mathcal{H}\!\mathit{om}}\nolimits}
\def\SheafExt{\mathop{\mathcal{E}\!\mathit{xt}}\nolimits}
\def\Sch{\mathit{Sch}}
\def\Mor{\mathop{\mathrm{Mor}}\nolimits}
\def\Ob{\mathop{\mathrm{Ob}}\nolimits}
\def\Sh{\mathop{\mathit{Sh}}\nolimits}
\def\NL{\mathop{N\!L}\nolimits}
\def\CH{\mathop{\mathrm{CH}}\nolimits}
\def\proetale{{pro\text{-}\acute{e}tale}}
\def\etale{{\acute{e}tale}}
\def\QCoh{\mathit{QCoh}}
\def\Ker{\mathop{\mathrm{Ker}}}
\def\Im{\mathop{\mathrm{Im}}}
\def\Coker{\mathop{\mathrm{Coker}}}
\def\Coim{\mathop{\mathrm{Coim}}}

% Boxtimes
%
\DeclareMathSymbol{\boxtimes}{\mathbin}{AMSa}{"02}

%
% Macros for moduli stacks/spaces
%
\def\QCohstack{\mathcal{QC}\!\mathit{oh}}
\def\Cohstack{\mathcal{C}\!\mathit{oh}}
\def\Spacesstack{\mathcal{S}\!\mathit{paces}}
\def\Quotfunctor{\mathrm{Quot}}
\def\Hilbfunctor{\mathrm{Hilb}}
\def\Curvesstack{\mathcal{C}\!\mathit{urves}}
\def\Polarizedstack{\mathcal{P}\!\mathit{olarized}}
\def\Complexesstack{\mathcal{C}\!\mathit{omplexes}}
% \Pic is the operator that assigns to X its picard group, usage \Pic(X)
% \Picardstack_{X/B} denotes the Picard stack of X over B
% \Picardfunctor_{X/B} denotes the Picard functor of X over B
\def\Pic{\mathop{\mathrm{Pic}}\nolimits}
\def\Picardstack{\mathcal{P}\!\mathit{ic}}
\def\Picardfunctor{\mathrm{Pic}}
\def\Deformationcategory{\mathcal{D}\!\mathit{ef}}

\setlength{\emergencystretch}{3em}
\begin{document}
\title[Stacks Project, Tag 0GTG]{473 --- Weak relative Noether normalization by a quasi-finite polynomial-algebra map with the maximal fibre dimension}
\maketitle
\noindent Copyright (C) 2005--2025 Johan de Jong. Stacks Project authors. Source: \href{https://stacks.math.columbia.edu/tag/0GTG}{Tag 0GTG}.

\noindent Editorial difficulty note (not author text). Difficulty H2: Approximation must preserve the global fibre-dimension bound while reducing to finite-type integer algebras. A double induction then constructs a coordinate dropping all fibre dimensions, first by transcendence at generic components and then by the perturbation a\textasciicircum{}n+b to reconcile an open stratum with its closed complement..

\noindent Editorial provenance note (not author text). History: excerpt from revision \texttt{a0444\allowbreak 6e57e\allowbreak c1fbc\allowbreak 252a8\allowbreak 71afc\allowbreak ec775\allowbreak 2fb28\allowbreak 07b14}, more-morphisms.tex, lines 8779--8882. Reference presentation was adapted; original-author context and core support blocks were retained or added. The mathematical wording is unchanged. Licensed under GNU FDL 1.2 or later; no invariant sections or cover texts. See ../licenses/COPYING. Context blocks reproduce author definitions and supporting results; they are not additional counted problems.

\begin{definition}
\label{algebra-definition-relative-dimension}
Suppose that $R \to S$ is of finite type, and let
$\mathfrak q \subset S$ be a prime lying over a prime
$\mathfrak p$ of $R$.
We define the {\it relative dimension
of $S/R$ at $\mathfrak q$}, denoted
$\dim_{\mathfrak q}(S/R)$, to be the dimension
of $\Spec(S \otimes_R \kappa(\mathfrak p))$
at the point corresponding to $\mathfrak q$. We let
$\dim(S/R)$ be the supremum of $\dim_{\mathfrak q}(S/R)$
over all $\mathfrak q$. This is called the
{\it relative dimension of} $S/R$.
\end{definition}

\begin{lemma}
\label{lemma-weak-relative-noether-normalization}
Let $R \to S$ be a finite type ring map. Let $d$ be the maximum
of the dimensions of fibres of $\Spec(S) \to \Spec(R)$.
Then there exists a quasi-finite ring map
$R[t_1, \ldots, t_d] \to S$.
\end{lemma}

\begin{proof}
In this paragraph we reduce to the case where $R \to S$ is of finite
presentation. Namely, write $S = R[x_1, \ldots, x_n]/J$ for some ideal
$J \subset R[x_1, \ldots, x_n]$.
Then $J$ is the union of its finitely generated ideals $J_\lambda \subset J$.
Set $S_\lambda = R[x_1, \ldots, x_n]/J_\lambda$.
Then for some $\lambda$ the fibres of
$\Spec(S_\lambda) \to \Spec(R)$
have dimension $\leq d$, see Limits, Lemma \href{https://stacks.math.columbia.edu/tag/05M5}{lemma limit dimension (Tag 05M5)}.
Fix such a $\lambda$. If we can find a quasi-finite
$R[t_1, \ldots, t_d] \to S_\lambda$, then of course the composition
$R[t_1, \ldots, t_d] \to S$ is quasi-finite.
This reduces us to the case discussed in the next paragraph.

\medskip\noindent
Assume $R \to S$ is of finite presentation. In this paragraph we reduce
to the case where $R$ is of finite type over $\mathbf{Z}$. By
Algebra, Lemma \href{https://stacks.math.columbia.edu/tag/00R1}{lemma limit module finite presentation (Tag 00R1)}
we can find a directed set $\Lambda$ and a system of ring maps
$R_\lambda \to S_\lambda$ over $\Lambda$ whose colimit is
$R \to S$ such that $S_\mu = S_\lambda \otimes_{R_\lambda} R_\mu$
for $\mu \geq \lambda$ and
such that each $R_\lambda$ and $S_\lambda$ is of finite
type over $\mathbf{Z}$.
Then for $\lambda$ large enough the fibres of
$\Spec(S_\lambda) \to \Spec(R_\lambda)$
have dimension $\leq d$, see
Limits, Lemma \href{https://stacks.math.columbia.edu/tag/0EY2}{lemma descend dimension d (Tag 0EY2)}.
Fix such a $\lambda$. If we can find a quasi-finite ring map
$R_\lambda[t_1, \ldots, t_d] \to S_\lambda$, then the base change
$R[t_1, \ldots, t_d] \to S$ is quasi-finite too
(Algebra, Lemma \href{https://stacks.math.columbia.edu/tag/00PP}{lemma quasi finite base change (Tag 00PP)}).
This reduces us the the case discussed in the next paragraph.

\medskip\noindent
Assume $R$ and $S$ are of finite type over $\mathbf{Z}$. If $d = 0$, then
the ring map is quasi-finite and we are done. Assume $d > 0$.
We will find an element $a \in S$ such that
the fibres of the $R$-algebra map $R[x] \to S$, $x \mapsto a$
have dimension $< d$. This will finish the proof by induction.

\medskip\noindent
We will prove the existence of $a$ by induction on $\dim(R)$.

\medskip\noindent
Let $\mathfrak q_1, \ldots, \mathfrak q_r \subset S$ be those
among the minimal primes of $S$ such that $\dim_{\mathfrak q_i}(S/R) = d$.
For notation, see
Algebra, Definition \ref{algebra-definition-relative-dimension}.
Say $\mathfrak q_i$ lies over the prime $\mathfrak p_i \subset R$.
We have $\text{trdeg}_{\kappa(\mathfrak p_i)}(\kappa(\mathfrak q_i)) = d$
as $\mathfrak q_i$ is a generic point of its fibre; for example
apply Algebra, Lemma \href{https://stacks.math.columbia.edu/tag/00P1}{lemma dimension at a point finite type field (Tag 00P1)}
to $S \otimes_R \kappa(\mathfrak p_i)$.
Hence by Lemma \href{https://stacks.math.columbia.edu/tag/0GTE}{lemma silly (Tag 0GTE)} we can find an element $a \in S$ such that
the image of $a$ in $\kappa(\mathfrak q_i)$ is transcendental over
$\kappa(\mathfrak p_i)$ for $i = 1, \ldots, r$.
Consider the morphism
$$
f_a : \Spec(S) \longrightarrow \mathbf{A}^1_R
$$
corresponding the $R$-algebra homomorphism $R[x] \to S$
to mapping $x$ to $a$. Let $U \subset \Spec(S)$ be the open
subset where the fibres have dimension $\leq d - 1$, see
Morphisms, Lemma \href{https://stacks.math.columbia.edu/tag/02FZ}{lemma openness bounded dimension fibres (Tag 02FZ)}.
By construction $U$ contains all the generic points of $\Spec(S)$.
In particular we see that $U$ contains all generic points of
all the generic fibres of $\Spec(S) \to \Spec(R)$ as such points are
necessarily generic points of $\Spec(S)$. Set $T = \Spec(S) \setminus U$
viewed as a reduced closed subscheme of $\Spec(S)$.
It follows from what we just said and the assumption that
$\dim(S/R) \leq d$ that
the generic fibres of $T \to \Spec(R)$ have dimension $\leq d - 1$.
Hence by Lemma \href{https://stacks.math.columbia.edu/tag/05F7}{lemma dimension in neighbourhood (Tag 05F7)},
applied several times to produce open neighbourhoods of the generic
points of $\Spec(R)$, we can find a dense open $V \subset \Spec(R)$ such that
$T_V \to V$ has fibres of dimension $\leq d - 1$.
We conclude that for $q \in \mathbf{A}^1_V$ the fibre of
$f_a$ over $q$ has dimension $\leq d - 1$ (as we have bounded the dimension
of the fibre of $f_a|_U$ and of the fibre of $f_a|_T$).

\medskip\noindent
By prime avoidance, we may assume that $V = D(f)$ for some $f \in R$.
Then we see that the ring map $R_f[x] \to S_f$, $x \mapsto a$
has fibres of dimension $\leq d - 1$.
We may replace $a$ by $fa$ and assume $a \in (f)$.
By induction on $\dim(R)$ we can find an element
$\overline{b} \in S/fS$
such that the fibres of $\Spec(S/fS) \to \Spec(R/fR[x])$,
$x \mapsto \overline{b}$ have dimension $\leq d - 1$.
Let $b \in S$ be a lift of $\overline{b}$.
By Lemma \href{https://stacks.math.columbia.edu/tag/0GTF}{lemma sum is ok (Tag 0GTF)} there exists an $n > 0$ such
that $a^n + b$ still works for $R_f \to S_f$.
On the other hand, the image of $a^n + b$ in $S/fS$ is $\overline{b}$
and the proof is complete.
\end{proof}
\end{document}
